\subsection{Downstream workflow evaluation on controlled gage basins}
\label{sec:methods-calibration}

Objective~4 tests whether generated SWAT\textsuperscript{+} projects can enter a standard downstream workflow---initialization, particle-swarm optimization \citep[PSO;][]{KennedyEberhart1995PSO} calibration, independent verification, and Morris global sensitivity \citep{Morris1991FactorialSampling, Campolongo2007ScreeningDesign}, following established split-sample practice \citep{Arnold2012SWATCalibration}---not whether the platform achieves national validity. Two compact HUC12 gage watersheds (\texttt{02297600} Florida; \texttt{05536265} Illinois) are evaluated separately from the S/M/L set.

A basin is \emph{workflow-ready} when the platform completes, without manual repair: (i)~initialization-pool generation and retention of the best draw on the calibration window; (ii)~PSO on a Morris-trimmed parameter subset (from \texttt{cal\_parms\_*.cal}) minimizing the negative sum of daily and monthly Nash--Sutcliffe efficiency \citep[NSE;][]{Moriasi2007ModelEval} at the gaged channel, through stagnation-based stop; (iii)~forward simulation of the global-best vector on the independent verification window; and (iv)~Morris batch execution with $\geq$999 successful runs, ensemble band diagnostics, and archived $\mu^*$ rankings. PSO used six concurrent forward runs on a ten-core host (36 particles for Florida, 48 for Illinois; both stopped after iteration~26); stations with $>$10\% missing daily observations in a scored window contribute objective zero.

Scored calibration and verification windows were chosen per basin from NWIS and meteorological coverage, with \textbf{no temporal overlap}, and calendar warm-up years precede each window via SWAT\textsuperscript{+} \texttt{nyskip}. Florida daily NWIS begins 2009-10-01, so the executed split is scored calibration 2013--2018 (warm-up from 2010) and post-calibration verification 2019--2024; Illinois uses scored calibration 2020--2024 (warm-up from 2018) and a pre-calibration holdout 2012--2015. Frozen run metadata are in the archived \texttt{tab-calibration-run-settings.csv}.

Reported metrics (daily/monthly NSE, KGE, PBIAS, RMSE by stage; Supplementary Table~\ref{tab:hydrologic-metrics} and Figures~S\ref{fig:supp-cal-val-hydro-fl}, S\ref{fig:supp-cal-val-hydro-il}) characterize controlled-basin plausibility only. Morris screening (1000 samples; 27 trajectories, four levels) ran on each basin's calibration window under a separate scenario; completed runs ($n=999$) form streamflow ensembles, where (per SWAT-CUP conventions) \textbf{P-factor} is the fraction of observations inside the 2.5--97.5\% band and \textbf{R-factor} is mean band width over observed standard deviation (Table~\ref{tab:sensitivity-ensemble-metrics}). Morris $\mu^*$ rankings (Tables~\ref{tab:morris-sensitivity}, \ref{tab:morris-sensitivity-05536265}) identify influential parameters but do not substitute for calibrating alternative preprocessing rule sets.
